%
% teil3.tex -- Beispiel-File für Teil 3
%
% (c) 2020 Prof Dr Andreas Müller, Hochschule Rapperswil
%
% !TEX root = ../../buch.tex
% !TEX encoding = UTF-8
%
\section{Zusammenfassung
  \label{kepler:section:zusammenfassung}}

Die Herleitung der keplerschen Gesetze demonstriert eindrücklich das Zusammenspiel von physikalischen Erhaltungssätzen und mathematischer Methodik.
Während die Betrachtung einer simplen Kreisbahn noch mit grundlegender elementarer Algebra auskommt, erfordert die allgemeine Lösung für die elliptische Bahn fortgeschrittenere mechanische Prinzipien.
Besonders die Anwendung komplexer Funktionen in Abschnitt \ref{kepler:subsection:komplexe_integration} zeigt, wie sich schwierige reelle physikalische Integrale durch den Übergang in die komplexe Ebene stark vereinfachen lassen.

